\documentclass[aspectratio=169]{beamer}
\usepackage[utf8]{inputenc}
\usepackage[T1]{fontenc}
\usepackage{amsmath,amssymb,amsfonts,bm,mathtools}
\usepackage{physics}
\usepackage{braket}
\usepackage{hyperref}
\usetheme{Madrid}

% ── convenience macros ────────────────────────────────────────────────────────
\newcommand{\E}{\mathbb{E}}

\newcommand{\sigmoid}{\sigma}
\newcommand{\pdata}{p_{\mathrm{data}}}
\newcommand{\pmodel}{p_{\theta}}
\newcommand{\Edata}{\langle\,\cdot\,\rangle_{\mathrm{data}}}
\newcommand{\Emodel}{\langle\,\cdot\,\rangle_{\mathrm{model}}}
\newcommand{\Ek}{\langle\,\cdot\,\rangle_{k}}
\newcommand{\DKL}{D_{\mathrm{KL}}}
\newcommand{\Zth}{Z_\theta}

\title{Contrastive Divergence}
\subtitle{Derivation from KL Divergence, Energy-Based Models, and Quantum Boltzmann Machines}
\author{Morten Hjorth-Jensen}
\date{April 30, 2026}

\begin{document}
\frame{\titlepage}

% ─────────────────────────────────────────────────────────────────────────────
\section{Motivation}
% ─────────────────────────────────────────────────────────────────────────────

\begin{frame}{Aim}
These notes derive contrastive divergence (CD) from the Kullback--Leibler
divergence and connect it to:
\begin{enumerate}
  \item classical Boltzmann machines (BMs),
  \item general energy-based models (EBMs),
  \item restricted Boltzmann machines (RBMs),
  \item quantum Boltzmann machines (QBMs).
\end{enumerate}

The central question is:
\[
  \boxed{\text{Why is the approximate BM gradient called
  \emph{contrastive divergence}?}}
\]
\end{frame}

% ─────────────────────────────────────────────────────────────────────────────
\section{Learning as KL Minimisation}
% ─────────────────────────────────────────────────────────────────────────────

\begin{frame}{The Learning Problem}
Let $\pdata(x)$ be the data distribution and $\pmodel(x)$ a parametric model.

\textbf{Goal:} find $\theta$ such that $\pmodel \approx \pdata$.

\medskip
Maximum likelihood (ML) maximises
\[
  \mathcal{L}(\theta)
  = \E_{\pdata}[\log \pmodel(x)]
  = \sum_x \pdata(x)\log\pmodel(x).
\]

\medskip
\textbf{Equivalence with KL minimisation.}
\[
  \DKL(\pdata\|\pmodel)
  = \sum_x \pdata(x)\log\frac{\pdata(x)}{\pmodel(x)}
  = \underbrace{H(\pdata)}_{\text{const.}} - \mathcal{L}(\theta),
\]
so minimising $\DKL(\pdata\|\pmodel)$ is exactly equivalent to maximising $\mathcal{L}(\theta)$.
\end{frame}

\begin{frame}{KL Divergence: Key Properties}
\begin{itemize}
  \item $\DKL(p\|q)\ge 0$, with equality iff $p=q$ (Gibbs inequality).
  \item $\DKL$ is \emph{not} symmetric: $\DKL(p\|q)\ne\DKL(q\|p)$ in general.
  \item $\DKL(p\|q)$ penalises regions where $p>0$ but $q\approx 0$
        (mode-covering behaviour for ML estimation).
\end{itemize}

\medskip
\textbf{Proof of non-negativity (Jensen).}
Since $-\log$ is convex:
\[
  \DKL(p\|q)
  = \E_p\!\left[-\log\frac{q(x)}{p(x)}\right]
  \ge -\log\E_p\!\left[\frac{q(x)}{p(x)}\right]
  = -\log\sum_x q(x) = 0.
\]
\end{frame}

% ─────────────────────────────────────────────────────────────────────────────
\section{Energy-Based Models}
% ─────────────────────────────────────────────────────────────────────────────

\begin{frame}{Energy-Based Models}
An \textbf{energy-based model} assigns an energy $E_\theta(x)\in\mathbb{R}$ to
each configuration $x$:
\[
  \pmodel(x)
  = \frac{e^{-E_\theta(x)}}{\Zth},
  \qquad
  \Zth = \sum_x e^{-E_\theta(x)}.
\]

Taking the log:
\[
  \log\pmodel(x) = -E_\theta(x) - \log\Zth.
\]

\textbf{The fundamental difficulty:} $\Zth$ is a sum over all configurations,
which is exponentially large for high-dimensional $x$.
\end{frame}

\begin{frame}{Log-Likelihood Gradient: Full Derivation}
\[
  \nabla_\theta\mathcal{L}(\theta)
  = \E_{\pdata}[\nabla_\theta\log\pmodel(x)]
  = -\E_{\pdata}[\nabla_\theta E_\theta(x)]
    - \nabla_\theta\log\Zth.
\]

\textbf{Gradient of the partition function.}
\begin{align*}
  \nabla_\theta\log\Zth
  &= \frac{\nabla_\theta\Zth}{\Zth}
   = \frac{\nabla_\theta\sum_x e^{-E_\theta(x)}}{\Zth} \\
  &= \frac{-\sum_x e^{-E_\theta(x)}\nabla_\theta E_\theta(x)}{\Zth}
   = -\E_{\pmodel}[\nabla_\theta E_\theta(x)].
\end{align*}

Substituting back:
\[
  \boxed{
    \nabla_\theta\mathcal{L}
    = -\E_{\pdata}[\nabla_\theta E_\theta]
      + \E_{\pmodel}[\nabla_\theta E_\theta].
  }
\]
\end{frame}

\begin{frame}{Positive and Negative Phases}
For \emph{minimising} the negative log-likelihood $-\mathcal{L}$:
\[
  \nabla_\theta(-\mathcal{L})
  = \underbrace{\langle\nabla_\theta E_\theta\rangle_{\mathrm{data}}}_{\text{positive phase}}
    -
    \underbrace{\langle\nabla_\theta E_\theta\rangle_{\mathrm{model}}}_{\text{negative phase}}.
\]

\begin{description}
  \item[Positive phase] Lowers $E_\theta$ on observed data configurations.
        Estimated cheaply from the training set.
  \item[Negative phase] Raises $E_\theta$ on configurations preferred by the
        current model. Requires samples from $\pmodel$ — generally intractable.
\end{description}

The negative phase is the computational bottleneck of all EBM training.
\end{frame}

% ─────────────────────────────────────────────────────────────────────────────
\section{Boltzmann Machines}
% ─────────────────────────────────────────────────────────────────────────────

\begin{frame}{Classical Boltzmann Machine}
A \textbf{Boltzmann machine} has visible units $v=(v_i)$ and hidden units
$h=(h_j)$, with energy
\[
  E_\theta(v,h)
  = -\sum_i a_i v_i
    - \sum_j b_j h_j
    - \sum_{ij} W_{ij}v_i h_j
    - \sum_{i<k} U_{ik}v_i v_k
    - \sum_{j<\ell} V_{j\ell}h_j h_\ell.
\]

The joint and visible distributions are
\[
  p_\theta(v,h)=\frac{e^{-E_\theta(v,h)}}{\Zth},
  \qquad
  p_\theta(v)=\sum_h p_\theta(v,h).
\]

The log-likelihood gradient over visible data is
\[
  \nabla_\theta\mathcal{L}
  = -\E_{\pdata(v)}\bigl[\E_{p(h|v)}[\nabla_\theta E_\theta(v,h)]\bigr]
    + \E_{p_\theta(v,h)}[\nabla_\theta E_\theta(v,h)].
\]
\end{frame}

\begin{frame}{Restricted Boltzmann Machine: Architecture}
In an \textbf{RBM} visible–visible and hidden–hidden connections are removed:
\[
  E_\theta(v,h)
  = -a^\top v - b^\top h - v^\top W h.
\]

This bipartite structure gives \emph{conditional independence}:
\[
  p_\theta(h\mid v) = \prod_j p_\theta(h_j\mid v),
  \qquad
  p_\theta(v\mid h) = \prod_i p_\theta(v_i\mid h).
\]

\textbf{Binary units.} For $v_i,h_j\in\{0,1\}$:
\[
  p(h_j=1\mid v)
  = \sigmoid\!\left(b_j + \textstyle\sum_i W_{ij}v_i\right),
\]
\[
  p(v_i=1\mid h)
  = \sigmoid\!\left(a_i + \textstyle\sum_j W_{ij}h_j\right),
\]
where $\sigmoid(x) = 1/(1+e^{-x})$.
\end{frame}

\begin{frame}{Deriving the Conditional Probabilities}
From the energy,
\[
  p(h_j=1\mid v)
  = \frac{p_\theta(v, h_j=1)}{p_\theta(v,h_j=0)+p_\theta(v,h_j=1)}.
\]

The numerator's energy contribution from $h_j$ is
$-(b_j + \sum_i W_{ij}v_i)h_j$.  Therefore

\[
  \frac{e^{b_j+\sum_i W_{ij}v_i}}{1 + e^{b_j+\sum_i W_{ij}v_i}}
  = \sigmoid\!\bigl(b_j + \textstyle\sum_i W_{ij}v_i\bigr).
\]

Because $v$--$h$ are bipartite, the $h_j$ are \emph{conditionally independent}
given $v$, so $p_\theta(h\mid v) = \prod_j p_\theta(h_j\mid v)$.
This makes block Gibbs sampling exact and efficient.
\end{frame}

\begin{frame}{Exact RBM Gradient}
Since $\partial E/\partial W_{ij} = -v_i h_j$, the gradient of the log-likelihood is
\[
  \frac{\partial\mathcal{L}}{\partial W_{ij}}
  = \langle v_i h_j\rangle_{\mathrm{data}}
    - \langle v_i h_j\rangle_{\mathrm{model}}.
\]

\textbf{Estimating the positive phase} (easy):
\[
  \langle v_i h_j\rangle_{\mathrm{data}}
  = \E_{\pdata(v)}\bigl[v_i\,\sigmoid(b_j+W_{:j}^\top v)\bigr].
\]

\textbf{Estimating the negative phase} (hard): requires equilibrium Gibbs
samples $\{v^*,h^*\}\sim p_\theta(v,h)$.

Similarly, for biases:
\[
  \frac{\partial\mathcal{L}}{\partial a_i}
  = \langle v_i\rangle_{\mathrm{data}} - \langle v_i\rangle_{\mathrm{model}},
  \qquad
  \frac{\partial\mathcal{L}}{\partial b_j}
  = \langle h_j\rangle_{\mathrm{data}} - \langle h_j\rangle_{\mathrm{model}}.
\]
\end{frame}

\begin{frame}{Gibbs Sampling for RBMs}
Block Gibbs sampling alternates:
\begin{enumerate}
  \item Sample $h \sim p_\theta(h\mid v)$: draw each $h_j$ independently
        from $\mathrm{Bernoulli}(\sigmoid(b_j+W_{:j}^\top v))$.
  \item Sample $v' \sim p_\theta(v\mid h)$: draw each $v_i$ independently
        from $\mathrm{Bernoulli}(\sigmoid(a_i+W_{i:}h))$.
\end{enumerate}

\medskip
Starting from any $v^{(0)}$, iterating $k$ steps gives samples from
$p_k(\cdot) = p_\theta(\cdot|T^k)$ where $T$ is the Gibbs transition.

\textbf{Convergence:} as $k\to\infty$, $p_k\to p_\theta$ (ergodic Markov chain).
CD-$k$ uses the $k$-step approximation in lieu of the true equilibrium.
\end{frame}

% ─────────────────────────────────────────────────────────────────────────────
\section{Contrastive Divergence}
% ─────────────────────────────────────────────────────────────────────────────

\begin{frame}{The Bottleneck and the CD Idea}
The exact model expectation
\[
  \langle v_i h_j\rangle_{\mathrm{model}}
\]
requires Gibbs chains run to equilibrium — expensive or impractical.

\medskip
\textbf{Contrastive divergence} (Hinton 2002) replaces equilibrium sampling by
\emph{short chains initialised at data}:

\[
  v^{(0)}\!\sim\pdata
  \;\to\; h^{(0)}\!\sim p(h|v^{(0)})
  \;\to\; v^{(1)}\!\sim p(v|h^{(0)})
  \;\to\; \cdots
  \;\to\; v^{(k)}.
\]

The distribution of $v^{(k)}$ is denoted $p_k$.  CD-$k$ uses
\[
  \langle v_i h_j\rangle_{\mathrm{model}}
  \approx \langle v_i h_j\rangle_{k}
  \;\Longrightarrow\;
  \Delta W_{ij}^{\mathrm{CD-}k}
  \propto
  \langle v_i h_j\rangle_{\mathrm{data}}
  -
  \langle v_i h_j\rangle_{k}.
\]
\end{frame}

\begin{frame}{Why ``Contrastive Divergence''? — The Objective}
Define the \textbf{contrastive divergence objective}:
\[
  \mathrm{CD}_k(\theta)
  = \DKL(p_0\|\pmodel)
    - \DKL(p_k\|\pmodel),
\]
where $p_0 = \pdata$ and $p_k = p_0 T_\theta^k$ ($T_\theta$ = Gibbs kernel leaving
$\pmodel$ invariant).

\medskip
\textbf{Interpretation.}
\begin{itemize}
  \item The first term pushes $\pmodel$ toward the data distribution.
  \item The second term \emph{undoes} some of that push, using
        the short-run reconstruction distribution $p_k$.
  \item The \emph{contrast} is between the data divergence and the $k$-step
        reconstruction divergence.
\end{itemize}

As $k\to\infty$, $p_k\to\pmodel$, so $\mathrm{CD}_k\to\DKL(p_0\|\pmodel)$ (exact ML).
\end{frame}

\begin{frame}{Gradient of the First KL Term}
Since $p_0 = \pdata$ does not depend on $\theta$:
\begin{align*}
  \nabla_\theta\DKL(p_0\|\pmodel)
  &= -\nabla_\theta\E_{p_0}[\log\pmodel(x)] \\
  &= -\E_{p_0}[\nabla_\theta\log\pmodel(x)] \\
  &= \E_{p_0}[\nabla_\theta E_\theta(x)]
     - \E_{\pmodel}[\nabla_\theta E_\theta(x)] \\
  &= \langle\nabla_\theta E_\theta\rangle_{p_0}
     - \langle\nabla_\theta E_\theta\rangle_{\pmodel}.
\end{align*}

This is the \emph{exact} ML gradient: positive phase minus negative phase.
\end{frame}

\begin{frame}{Gradient of the Second KL Term}
\[
  \DKL(p_k\|\pmodel)
  = \sum_x p_k(x)\log\frac{p_k(x)}{\pmodel(x)}.
\]

The full gradient has two contributions:
\begin{enumerate}
  \item \textbf{Explicit:} derivative through $\pmodel$ (treating $p_k$ as fixed).
  \item \textbf{Implicit:} derivative through $p_k = p_0 T_\theta^k$, since
        $T_\theta$ depends on $\theta$.
\end{enumerate}

The \textbf{CD approximation} neglects the implicit term:
\[
  \nabla_\theta p_k \approx 0.
\]

Under this approximation:
\[
  \nabla_\theta\DKL(p_k\|\pmodel)
  \approx
  \langle\nabla_\theta E_\theta\rangle_{p_k}
  - \langle\nabla_\theta E_\theta\rangle_{\pmodel}.
\]
\end{frame}

\begin{frame}{The CD Gradient: Cancellation}
\begin{align*}
  \nabla_\theta\mathrm{CD}_k
  &\approx
  \nabla_\theta\DKL(p_0\|\pmodel)
  -
  \nabla_\theta\DKL(p_k\|\pmodel) \\[4pt]
  &=
  \Bigl[\langle\nabla_\theta E_\theta\rangle_{p_0}
         - \langle\nabla_\theta E_\theta\rangle_{\pmodel}\Bigr]
  -
  \Bigl[\langle\nabla_\theta E_\theta\rangle_{p_k}
         - \langle\nabla_\theta E_\theta\rangle_{\pmodel}\Bigr].
\end{align*}

The model expectations $\langle\nabla_\theta E_\theta\rangle_{\pmodel}$ cancel:
\[
  \boxed{
    \nabla_\theta\mathrm{CD}_k
    \approx
    \langle\nabla_\theta E_\theta\rangle_{p_0}
    -
    \langle\nabla_\theta E_\theta\rangle_{p_k}.
  }
\]

This is the final CD update rule:\ data statistics minus $k$-step
reconstruction statistics.  The intractable partition-function gradient
has been eliminated.
\end{frame}

\begin{frame}{Concrete RBM CD-1 Update}
For CD-1 on an RBM, one Gibbs step gives $v^{(1)}$ from $v^{(0)}\sim\pdata$:

\[
  h_j^{(0)} \sim \mathrm{Bernoulli}\!\bigl(\sigmoid(b_j + W_{:j}^\top v^{(0)})\bigr),
\]
\[
  v_i^{(1)} \sim \mathrm{Bernoulli}\!\bigl(\sigmoid(a_i + W_{i:} h^{(0)})\bigr),
\]
\[
  h_j^{(1)} \sim \mathrm{Bernoulli}\!\bigl(\sigmoid(b_j + W_{:j}^\top v^{(1)})\bigr).
\]

The CD-1 weight update is then:
\begin{align*}
  \Delta W_{ij}
  &= \eta\Bigl(v_i^{(0)}\,\E[h_j^{(0)}\mid v^{(0)}]
               - v_i^{(1)}\,\E[h_j^{(1)}\mid v^{(1)}]\Bigr), \\
  \Delta a_i
  &= \eta\bigl(v_i^{(0)} - v_i^{(1)}\bigr), \qquad
  \Delta b_j
   = \eta\bigl(\E[h_j^{(0)}\mid v^{(0)}] - \E[h_j^{(1)}\mid v^{(1)}]\bigr).
\end{align*}

Using the \emph{mean-field} hidden activations (rather than samples) reduces variance.
\end{frame}

% ─────────────────────────────────────────────────────────────────────────────
\section{Bias, Variants, and Connections}
% ─────────────────────────────────────────────────────────────────────────────

\begin{frame}{Bias of the CD Gradient}
CD-$k$ is \emph{not} an unbiased estimator of the ML gradient.

\medskip
\textbf{The bias.} Let $T_\theta$ be the Gibbs transition kernel.  The true
gradient of $\DKL(p_k\|\pmodel)$ is
\[
  \nabla_\theta\DKL(p_k\|\pmodel)
  = \langle\nabla_\theta E_\theta\rangle_{p_k}
    - \langle\nabla_\theta E_\theta\rangle_{\pmodel}
    + \underbrace{\sum_x (\nabla_\theta p_k(x))
      \log\frac{p_k(x)}{\pmodel(x)}}_{\text{neglected implicit term}}.
\]

Neglecting the implicit term introduces a bias that vanishes as $k\to\infty$.

\medskip
\textbf{Consequence.} CD-$k$ may not converge to the true ML solution and can
exhibit slow drift or oscillation during training.
\end{frame}

\begin{frame}{Persistent Contrastive Divergence (PCD)}
\textbf{Key idea} (Tieleman 2008): maintain a persistent set of fantasy particles
$\{v^{(n)}\}_{n=1}^N$ that are \emph{not} reinitialised from data.

\medskip
\textbf{PCD update:}
\begin{enumerate}
  \item Run one Gibbs step from each fantasy particle: $v^{(n)}\to v^{(n)}{}'$.
  \item Compute:
        $\Delta W_{ij} = \eta\bigl(\langle v_i h_j\rangle_{\mathrm{data}}
                         - \langle v_i h_j\rangle_{\mathrm{fantasy}}\bigr)$.
  \item Update the fantasy particles: $v^{(n)}\leftarrow v^{(n)}{}'$.
\end{enumerate}

\medskip
\textbf{Why better?} Fantasy particles mix across modes between updates; they
approximate the true model distribution more faithfully than CD-$k$ short chains.

\medskip
\textbf{CD vs PCD:}
\begin{tabular}{lll}
  & \textbf{CD-$k$} & \textbf{PCD} \\
  Init.\ of chain & data & persistent fantasy \\
  Bias & larger & smaller \\
  Mixing & local & global \\
\end{tabular}
\end{frame}

\begin{frame}{Stochastic Maximum Likelihood (SML)}
SML is essentially PCD with a rigorous motivation.

\medskip
At step $t$, the Markov chain is at $v^{(t)}$, which approximates a sample from
$p_{\theta^{(t)}}$.  After a small parameter update $\theta^{(t)}\to\theta^{(t+1)}$,
$v^{(t)}$ is still close to a sample from $p_{\theta^{(t+1)}}$ provided the
update is small.

\medskip
\textbf{Gradient approximation:}
\[
  \nabla_\theta\mathcal{L}(\theta)
  \approx
  \langle\nabla_\theta E_\theta\rangle_{\mathrm{data}}
  -
  \nabla_\theta E_\theta(v^{(t)}).
\]

The approximation improves as the learning rate decreases, and under mild
conditions SML is asymptotically unbiased.
\end{frame}

\begin{frame}{Natural Gradient and Fisher Information}
Standard gradient ascent on $\mathcal{L}(\theta)$ is geometry-unaware.

\medskip
The \textbf{natural gradient} (Amari 1998) uses the \emph{Fisher information
matrix} $F(\theta)$:
\[
  F_{ab}(\theta)
  = \E_{\pmodel}\!\left[
      \frac{\partial\log\pmodel}{\partial\theta_a}
      \frac{\partial\log\pmodel}{\partial\theta_b}
    \right],
\]
and updates
\[
  \theta \leftarrow \theta + \eta\, F(\theta)^{-1}\nabla_\theta\mathcal{L}(\theta).
\]

\textbf{For EBMs:}
\[
  F_{ab}(\theta)
  = \operatorname{Cov}_{\pmodel}\!\left[
      \nabla_{\theta_a}E_\theta,\;
      \nabla_{\theta_b}E_\theta
    \right].
\]

Computing $F^{-1}$ is expensive, but the natural gradient converges faster and is
invariant to reparameterisation.  Approximate versions (e.g., K-FAC) are used in
practice.
\end{frame}

% ─────────────────────────────────────────────────────────────────────────────
\section{Energy-Based Interpretation}
% ─────────────────────────────────────────────────────────────────────────────

\begin{frame}{Energy Landscape Picture}
For $\pmodel(x)\propto e^{-E_\theta(x)}$, learning reshapes the energy surface:

\medskip
\begin{center}
\begin{tabular}{ll}
  \textbf{ML} & Lower $E$ on data; raise $E$ on model samples. \\[4pt]
  \textbf{CD-$k$} & Lower $E$ on data; raise $E$ on $k$-step reconstructions.
\end{tabular}
\end{center}

\medskip
The CD update is \emph{local}: it adjusts the energy landscape only in the
neighbourhood of data points, as defined by the mixing distance of $k$ Gibbs steps.

\medskip
This locality is both a strength (cheap) and a weakness (misses distant modes).
\end{frame}

\begin{frame}{Relation to Other EBM Training Methods}
\begin{tabular}{lll}
  \textbf{Method} & \textbf{Negative samples} & \textbf{Notes} \\[3pt]
  CD-$k$ & short Gibbs chains & biased, cheap \\
  PCD / SML & persistent chains & less biased \\
  Score matching & analytic & no sampling needed \\
  Noise-contrastive est. & fixed noise dist.\ & binary classifier \\
  Diffusion-based & learned reverse SDE & SOTA generative models \\
\end{tabular}

\medskip
CD uses the model's own dynamics to generate negative samples.  Unlike score
matching, it does not require second-order derivatives of $E_\theta$.  Unlike
NCE, it does not need an explicit noise distribution.
\end{frame}

% ─────────────────────────────────────────────────────────────────────────────
\section{Quantum Boltzmann Machines}
% ─────────────────────────────────────────────────────────────────────────────

\begin{frame}{Quantum Boltzmann Machine: Setup}
A \textbf{quantum Boltzmann machine} replaces the classical energy by a
Hamiltonian $H_\theta$ acting on a Hilbert space $\mathcal{H}$.

\medskip
The \textbf{Gibbs (thermal) state} at inverse temperature $\beta$ is
\[
  \rho_\theta
  = \frac{e^{-\beta H_\theta}}{\Zth},
  \qquad
  \Zth = \Tr\!\bigl(e^{-\beta H_\theta}\bigr).
\]

Visible probabilities arise from projective measurements:
\[
  p_\theta(v)
  = \Tr\!\bigl[\Pi_v\,\rho_\theta\bigr],
\]
where $\Pi_v = |v\rangle\langle v|$ projects onto visible bitstring $v$, and
the hidden degrees of freedom have been traced out.
\end{frame}

\begin{frame}{Quantum Log-Likelihood and Its Gradient}
The training objective is
\[
  \mathcal{L}(\theta)
  = \sum_v \pdata(v)\log p_\theta(v)
  = \sum_v \pdata(v)\log\Tr\!\bigl[\Pi_v\rho_\theta\bigr].
\]

The gradient is
\[
  \nabla_\theta\mathcal{L}
  = \sum_v \pdata(v)
    \frac{\Tr[\Pi_v\,\partial_\theta\rho_\theta]}{\Tr[\Pi_v\,\rho_\theta]}.
\]

This requires computing $\partial_\theta\rho_\theta$, which is where the
commuting and non-commuting cases diverge.
\end{frame}

\begin{frame}{Quantum Gradient: Commuting Hamiltonian}
Suppose $H_\theta = \sum_a \theta_a O_a$ with all $[O_a, O_{a'}]=0$.

\medskip
Then $e^{-\beta H_\theta} = \prod_a e^{-\beta\theta_a O_a}$, and
\[
  \partial_{\theta_a}e^{-\beta H_\theta}
  = -\beta O_a\, e^{-\beta H_\theta}.
\]

Therefore
\[
  \partial_{\theta_a}\Zth
  = -\beta\Tr\!\bigl(O_a e^{-\beta H_\theta}\bigr)
  = -\beta\Zth\langle O_a\rangle_{\rho_\theta},
\]
and
\[
  \partial_{\theta_a}\log\Zth = -\beta\langle O_a\rangle_{\rho_\theta}.
\]

The gradient of the log-likelihood has the classical BM structure:
\[
  \nabla_{\theta_a}(-\mathcal{L})
  = \beta\bigl(\langle O_a\rangle_{\mathrm{data}}
               - \langle O_a\rangle_{\rho_\theta}\bigr).
\]
\end{frame}

\begin{frame}{Quantum Gradient: Non-Commuting Hamiltonian}
When $[H_\theta, \partial_{\theta_a}H_\theta]\ne 0$, the derivative of
$e^{-\beta H_\theta}$ is given by the \textbf{Duhamel (Dyson) formula}:
\[
  \partial_{\theta_a}e^{-\beta H_\theta}
  =
  -\int_0^\beta
    e^{-(\beta-\tau)H_\theta}
    (\partial_{\theta_a}H_\theta)
    e^{-\tau H_\theta}
  \,d\tau.
\]

Taking the trace:
\[
  \partial_{\theta_a}\Zth
  =
  -\int_0^\beta
    \Tr\!\Bigl[
      e^{-(\beta-\tau)H_\theta}
      O_a\,
      e^{-\tau H_\theta}
    \Bigr]
  d\tau
  =
  -\Zth
  \int_0^\beta
    \langle O_a(\tau)\rangle_\beta
  \,d\tau,
\]
where $O_a(\tau) = e^{\tau H_\theta}O_a e^{-\tau H_\theta}$ is the
\emph{imaginary-time Heisenberg picture} operator and
$\langle\cdot\rangle_\beta = \Tr[\rho_\theta\,\cdot\,]/\Tr[\rho_\theta]$.
\end{frame}

\begin{frame}{Imaginary-Time Correlation Functions}
Defining the \textbf{symmetrised imaginary-time correlator}
\[
  \widetilde{\langle O_a\rangle}_{\rho_\theta}
  \;:=\;
  \frac{1}{\beta}\int_0^\beta
    \langle O_a(\tau)\rangle_\beta
  \,d\tau
  \;=\;
  \frac{1}{\beta}
  \int_0^\beta
    \Tr\!\bigl[\rho_\theta^{1-\tau/\beta} O_a \rho_\theta^{\tau/\beta}\bigr]
  \,d\tau,
\]
the partition-function gradient becomes
\[
  \partial_{\theta_a}\log\Zth
  = -\beta\,\widetilde{\langle O_a\rangle}_{\rho_\theta}.
\]

\textbf{Key point.} In the commuting case,
$\widetilde{\langle O_a\rangle} = \langle O_a\rangle_{\rho_\theta}$.
In the non-commuting case, the gradient requires imaginary-time
correlation functions, which are generally inaccessible from projective
measurements alone.
\end{frame}

\begin{frame}{Full QBM Gradient Formula}
Combining, for a QBM with non-commuting Hamiltonian:
\[
  \nabla_{\theta_a}(-\mathcal{L})
  =
  \beta\!\sum_v
    \pdata(v)
    \frac{\displaystyle
      \Tr\!\Bigl[\Pi_v
      \int_0^\beta e^{-(\beta-\tau)H_\theta}O_a e^{-\tau H_\theta}\,d\tau
      \Bigr]}
    {\Tr[\Pi_v\rho_\theta]}
  \,-\,
  \beta\,\widetilde{\langle O_a\rangle}_{\rho_\theta}.
\]

\textbf{Two difficulties:}
\begin{itemize}
  \item The positive phase requires conditional imaginary-time correlators
        given each visible configuration $v$ — exponentially hard classically.
  \item The negative phase requires the unconditional imaginary-time correlator
        $\widetilde{\langle O_a\rangle}_{\rho_\theta}$.
\end{itemize}

Both can be accessed on a quantum processor but require care about
state preparation and measurement.
\end{frame}

\begin{frame}{Quantum Contrastive Divergence}
\textbf{Classical analogy:}
\[
  \pdata \;\longrightarrow\; p_k
  \quad\text{via}\quad
  k \text{ steps of Gibbs sampling}.
\]

\textbf{Quantum analogue:}
\[
  \rho_{\mathrm{data}} \;\longrightarrow\; \rho_k,
\]
where $\rho_k$ is obtained by:
\begin{itemize}
  \item \emph{short quantum annealing} toward the Gibbs state,
  \item \emph{variational thermal-state} preparation (e.g., VQT),
  \item \emph{Suzuki--Trotter} imaginary-time evolution for $k$ steps,
  \item or \emph{mean-field} quantum approximations.
\end{itemize}

Under the same approximation $\nabla_\theta\rho_k\approx 0$:
\[
  \Delta\theta_a
  \propto
  \langle O_a\rangle_{\mathrm{data}}
  -
  \langle O_a\rangle_{\rho_k}.
\]
\end{frame}

\begin{frame}{Quantum Contrastive Divergence: Objective}
Defining the \textbf{quantum CD objective} as
\[
  \mathrm{QCD}_k(\theta)
  = S(\rho_0\|\rho_\theta)
    - S(\rho_k\|\rho_\theta),
\]
where $S(\rho\|\sigma) = \Tr[\rho(\log\rho - \log\sigma)]$ is the
\textbf{quantum relative entropy} (KL divergence for density matrices).

\medskip
This mirrors the classical construction exactly, with
$\DKL(p\|q)$ replaced by $S(\rho\|\sigma)$.

\medskip
\textbf{Properties of $S(\rho\|\sigma)$:}
\begin{itemize}
  \item $S(\rho\|\sigma)\ge 0$, with equality iff $\rho=\sigma$ (Klein's inequality).
  \item $S(\rho\|\sigma)=\Tr[\rho\log\rho]-\Tr[\rho\log\sigma]$.
  \item Not symmetric in general.
\end{itemize}
\end{frame}

\begin{frame}{Practical Challenges for QBMs}
\begin{tabular}{lp{8cm}}
  \textbf{Challenge} & \textbf{Details} \\[3pt]
  Gibbs-state prep.\ & Exact preparation is QMA-hard in general;
                        variational or annealing-based methods are needed. \\[3pt]
  Non-commuting $H$ & Requires imaginary-time correlators,
                       not accessible from standard measurements. \\[3pt]
  Hidden variables & Partial trace over hidden Hilbert space;
                     exponential in hidden dimension. \\[3pt]
  Hardware noise & Decoherence distorts the thermal state;
                    error mitigation is essential. \\[3pt]
  Gradient estimation & Each gradient call may require many circuit executions.
\end{tabular}

\medskip
\textbf{Practical approaches:}
mean-field QBMs, Suzuki--Trotter thermalization,
variational thermal-state preparation (VQT),
quantum annealing samples,
and CD-like short-run estimates.
\end{frame}

% ─────────────────────────────────────────────────────────────────────────────
\section{Summary}
% ─────────────────────────────────────────────────────────────────────────────

\begin{frame}{Summary: Classical CD}
For a classical EBM $\pmodel(x) = e^{-E_\theta(x)}/\Zth$,
ML maximisation gives the intractable gradient
\[
  \nabla_\theta\mathcal{L}
  = -\langle\nabla_\theta E_\theta\rangle_{\mathrm{data}}
    + \langle\nabla_\theta E_\theta\rangle_{\mathrm{model}}.
\]

CD-$k$ approximates this by replacing equilibrium model expectations
with $k$-step Gibbs reconstructions:
\[
  \nabla_\theta\mathrm{CD}_k
  \approx
  \langle\nabla_\theta E_\theta\rangle_{p_0}
  -
  \langle\nabla_\theta E_\theta\rangle_{p_k}.
\]

This approximation is \textbf{biased} (neglects $\nabla_\theta p_k$) but
\textbf{efficient} (no partition function needed).
\end{frame}

\begin{frame}{Summary: Why the Name?}
The CD objective is
\[
  \mathrm{CD}_k(\theta)
  =
  \underbrace{\DKL(p_0\|\pmodel)}_{\text{data divergence}}
  -
  \underbrace{\DKL(p_k\|\pmodel)}_{\text{reconstruction divergence}}.
\]

\begin{description}
  \item[Contrastive] Compare (contrast) data statistics with short-run
        reconstruction statistics.
  \item[Divergence] The comparison is formalised through KL divergences.
\end{description}

As $k$ increases the approximation improves:
\[
  k=0: \quad \mathrm{CD}_0 = 0;
  \qquad
  k\to\infty: \quad \mathrm{CD}_k \to \DKL(p_0\|\pmodel) \;(= \text{exact ML}).
\]
\end{frame}

\begin{frame}{Summary: Quantum Extension}
For a QBM $\rho_\theta = e^{-\beta H_\theta}/\Zth$, the key results are:

\medskip
\begin{tabular}{ll}
  Commuting $H$ & $\nabla_{\theta_a}(-\mathcal{L})
                  = \beta(\langle O_a\rangle_{\mathrm{data}}
                         - \langle O_a\rangle_{\rho_\theta})$ \\[6pt]
  Non-commuting $H$ & $\nabla_{\theta_a}\log\Zth
                       = -\beta\,\widetilde{\langle O_a\rangle}_{\rho_\theta}$,
                       imaginary-time integral \\[6pt]
  Quantum CD & $\Delta\theta_a
               \propto \langle O_a\rangle_{\mathrm{data}}
                       - \langle O_a\rangle_{\rho_k}$
\end{tabular}

\medskip
The quantum CD objective replaces KL divergence with quantum relative entropy:
\[
  \mathrm{QCD}_k(\theta)
  = S(\rho_0\|\rho_\theta) - S(\rho_k\|\rho_\theta).
\]
\end{frame}

\begin{frame}{Final Take-Home Message}
\[
  \boxed{
    \mathrm{CD}_k
    = \DKL(p_{\mathrm{data}}\|\pmodel)
      - \DKL(p_k\|\pmodel)
  }
\]

The learning rule contrasts
\[
  \text{data-driven expectations}
  \quad\text{with}\quad
  \text{short-run reconstructed expectations},
\]
eliminating the need to evaluate the partition function.

\medskip
This single idea unifies:
\begin{enumerate}
  \item Classical EBMs and BMs (Gibbs sampling approximation),
  \item RBMs (efficient block Gibbs with factored conditionals),
  \item QBMs (quantum relative entropy, Duhamel formula, imaginary-time correlators),
  \item Persistent CD / SML (global approximation to the model distribution).
\end{enumerate}
\end{frame}

\end{document}
